% ===========================================================================
% BMIN 5200 syllabus — Fall 2026.
%
% THIS FILE IS THE SOURCE OF TRUTH for the syllabus. It replaced
% syllabus-2026.md, which was deleted so there is only one prose copy to keep
% current. When the syllabus changes, edit here and rebuild.
%
% BUILD (requires a local TeX Live; nothing in CI compiles this):
%
%     latexmk -pdf -outdir=build-tex syllabus-2026.tex
%     cp build-tex/syllabus-2026.pdf new-course-site/syllabus-2026.pdf
%
% Then COMMIT THE PDF. new-course-site/syllabus-2026.pdf is a static asset
% served straight off the site; apply_links.py copies it into build/ via
% PASSTHROUGH. There is deliberately no LaTeX toolchain in the Actions
% workflow -- a TeX error must not be able to take the whole site down.
%
% Stock TeX Live packages only. No custom class, no minted, no shell-escape.
% ===========================================================================

\documentclass[11pt]{article}

\usepackage[T1]{fontenc}
\usepackage[utf8]{inputenc}
\usepackage{lmodern}
\usepackage[english]{babel}
\usepackage[margin=1in]{geometry}
\usepackage{booktabs}
\usepackage{longtable}
\usepackage{array}
\usepackage{enumitem}
\usepackage{titlesec}
\usepackage{xcolor}
\usepackage[hidelinks]{hyperref}

% Match the course site's accent (--accent: #8c1515).
\definecolor{accent}{HTML}{8C1515}
\definecolor{inksoft}{HTML}{5A5F66}

\hypersetup{
  colorlinks=true,
  linkcolor=accent,
  urlcolor=accent,
  pdftitle={BMIN 5200: Foundations of Artificial Intelligence in Health --- Fall 2026},
  pdfauthor={Joseph D. Romano},
}

\titleformat{\section}{\Large\bfseries\color{accent}}{}{0pt}{}
\titlespacing*{\section}{0pt}{1.4em}{0.5em}
\titleformat{\subsection}{\large\bfseries}{}{0pt}{}
\titlespacing*{\subsection}{0pt}{1.1em}{0.4em}

\setlist[itemize]{leftmargin=1.2em, itemsep=0.2em, topsep=0.3em}
\setlist[enumerate]{leftmargin=1.6em, itemsep=0.2em, topsep=0.3em}

\renewcommand{\arraystretch}{1.25}
\setlength{\parindent}{0pt}
\setlength{\parskip}{0.55em}

\newcommand{\handbookurl}{https://www.med.upenn.edu/BMI-training-programs/student-and-faculty-handbook.pdf}

\begin{document}

% ---------------------------------------------------------------- title ----
\begin{center}
  {\color{inksoft}\small\scshape BMIN 5200-001 \quad$\cdot$\quad Fall 2026}\\[0.6em]
  {\LARGE\bfseries Foundations of Artificial Intelligence in Health}\\[0.9em]
  {\normalsize Thursdays, 3:30--6:30pm \quad$\cdot$\quad 3600 Civic Center Blvd, Room 6E 031}\\[0.2em]
  {\normalsize August 27 -- December 3, 2026}\\[0.2em]
  {\normalsize Joseph D. Romano, PhD \quad$\cdot$\quad
   \href{mailto:joseph.romano@pennmedicine.upenn.edu}{joseph.romano@pennmedicine.upenn.edu}}
\end{center}

\vspace{0.4em}
{\color{accent}\hrule height 1.5pt}
\vspace{0.8em}

Welcome to BMIN 5200. This syllabus outlines the logistical details of the course. All course
materials live on the \href{https://bmin5200.jdr.bio}{course website}; Canvas is used only for
submitting assignments, viewing grades, and watching class recordings.

Don't hesitate to ask if you have any questions about the information here or elsewhere --- Joe
can be reached at
\href{mailto:joseph.romano@pennmedicine.upenn.edu}{joseph.romano@pennmedicine.upenn.edu}, or you
can ask before or after class.

% ---------------------------------------------------------- description ----
\section{Course description}

As a subfield of computer science, artificial intelligence (AI) is often used interchangeably
with the term `machine learning', which is itself only a subfield of AI dealing with the broader
concept of \emph{inductive reasoning}. However, a wealth of key prerequisite topics that focus on
\emph{deductive reasoning} cover a large portion of biomedical informatics applications being
actively used today. These founding principles of AI and their intersection with biomedical
informatics are the focus of this first course on AI, and lay the groundwork for future
coursework on machine learning, artificial neural networks, and generative AI.

This course is divided into modules that cover (1) introductory/background materials,
(2) knowledge representation, (3) logic, (4) essentials of rule-based systems, (5) search,
(6) information structure and inference, and (7) special topics. These topics offer a global
foundation of branches of AI application and research in biomedical domains, including concepts
that will later support a deeper understanding of inductive reasoning and machine learning.

In a practical sense, this course focuses on how biomedical data can be organized, represented,
interpreted, searched, and applied in order to derive knowledge, make decisions, and ultimately
make predictions while mitigating bias.

\emph{To reiterate: this course does not focus heavily on machine learning or generative AI, but
it does present critically important deductive and symbolic AI concepts that are necessary for a
deeper understanding of ML and GenAI. BMIN 5200 is the first in a 3-part course series on
artificial intelligence, followed by BMIN 5210 (AI 2: Machine Learning) and BMIN 5220 (AI 3:
Natural Language Processing).}

% ----------------------------------------------------------- objectives ----
\section{Course objectives}

Upon completion of this course, students will be able to:

\begin{itemize}
  \item Describe the unique challenges of representing and applying human knowledge in an
        \emph{in silico} environment
  \item Explain the conceptual differences between deduction and induction, and the role of both
        in AI
  \item Discuss the past and current challenges to developing and implementing AI in biomedical
        domains
  \item Demonstrate knowledge of how common AI approaches can be developed to address real-world
        biomedical problems
  \item \textbf{Implement} core AI methods in Python --- inference engines, search algorithms,
        expert systems, and probabilistic models --- rather than only describing them
\end{itemize}

% -------------------------------------------------------- prerequisites ----
\section{Prerequisites}

It is expected that all students will be somewhat familiar with basic biomedical concepts and
terminology, as well as statistics and computer programming. It is suggested (but not strictly
required) that students have taken Introduction to Biomedical Informatics (BMIN 5010), Data
Science for Biomedical Informatics (BMIN 5030), and a programming course (any language is fine).
All programming in this course will be done in Python, and a brief introduction to the language
will be provided.

No previous exposure to artificial intelligence is assumed.

% ------------------------------------------------------- class meetings ----
\section{Class meetings}

\textbf{Thursdays, 3:30--6:30pm}, in person.\\
\textbf{3600 Civic Center Boulevard, Room 6E 031.}

First meeting \textbf{August 27, 2026}. Last meeting \textbf{December 3, 2026}.

Class will not meet on:

\begin{itemize}
  \item \textbf{October 1, 2026} --- Fall Break (October 1--4)
  \item \textbf{November 26, 2026} --- Thanksgiving (November 26--29)
\end{itemize}

That leaves 13 meetings, one per lecture topic.

\subsection{Recording and virtual attendance}

Class sessions are recorded over Zoom and posted to Canvas, so you can review a lecture later or
catch up on one you missed. The recordings are for enrolled students only --- do not redistribute
them.

\textbf{There is no option to attend virtually.} The recordings exist for review, not as a
substitute for being in the room: this course is built around discussion, journal club, and
in-class work, none of which function if you are watching a video afterward. Attendance is in
person.

\subsection{Format of a typical meeting}

\begin{tabular}{@{}ll@{}}
  \toprule
  0--30 min & Homework review \\
  $\sim$60 min & Weekly lecture \\
  $\sim$25 min & \textbf{In-class coding exercise} \\
  10 min & Break \\
  $\sim$25 min & Journal club \\
  \bottomrule
\end{tabular}

The in-class coding exercise is new this year, and it replaces the second journal club paper.
Each week's exercise is a Jupyter notebook that runs in Google Colab --- you click a link, and it
works. You will build a propositional inference engine, a forward and backward chaining rule
system, a search algorithm suite, an expert system in CLIPS, a Bayesian network, and a clinical
agent, among others. These exercises earn credit for engagement and completion rather than for
correct answers, and they are where the course's ideas stop being abstract. Several of them are
direct on-ramps to the homework.

\textbf{You must bring a laptop to every class.} You do not need to install anything: the
exercises run in a browser through Colab. If you would rather work locally, a
\texttt{requirements.txt} is provided in the course repository.

% ---------------------------------------------------------------- units ----
\section{Course units}

This is a 1.0cu course.

% ------------------------------------------------------------ materials ----
\section{Where course materials live}

\textbf{Course website:} the schedule, links to every lecture deck, the in-class notebooks, the
journal club papers, and the assignment specifications. This is the first place to look.

\textbf{Penn Box:} lecture slides, journal club PDFs, and assignment files. Box links from the
course website require Penn SSO --- sign in with your PennKey.

\textbf{Canvas:} assignment submission, grades, and the Zoom recordings of each class. Nothing
else in the course depends on it --- lecture materials, notebooks, and papers are all reached
from the course website.

\textbf{GitHub:} the in-class exercise notebooks, in a public repository, each with an ``Open in
Colab'' link. The repository also holds
\href{https://github.com/RomanoLab/bmin5200/blob/main/SETUP.md}{SETUP.md}, which covers Colab,
homework submission, and the optional local Python environment. Nothing there needs a PennKey.

% --------------------------------------------------------- office hours ----
\section{Office hours}

TBD --- a poll will go out in the first week, and the time that works for the greatest number of
students will be selected. One-on-one meetings with the course director or TA may be scheduled on
an as-needed basis.

% ------------------------------------------------------ course director ----
\section{Course director}

\textbf{Joseph D. Romano, PhD, MPhil, MA} (please feel free to call me Joe)\\
Assistant Professor of Informatics and Pharmacology\\
3600 Civic Center Blvd, 5E 300A, Philadelphia, PA 19104\\
(+1) 215-573-5571\\
\href{mailto:joseph.romano@pennmedicine.upenn.edu}{joseph.romano@pennmedicine.upenn.edu}

\section{Teaching assistant}

TBD --- will be announced before the first class meeting.

% ------------------------------------------------------- expectations ------
\section{Student expectations}

Students are expected to come to class prepared for the day's content, and to ask questions and
discuss that content following the lectures. Reviewing lecture slides before coming to class is
strongly recommended.

Students are expected to read all assigned materials, participate during class sessions, and
complete the required assignments. This course requires the use of a laptop computer.

% ----------------------------------------------------------- evaluation ----
\section{Student evaluation}

\begin{tabular}{@{}lr@{}}
  \toprule
  \textbf{Assignments} (four, weighted equally) & \textbf{40\%} \\
  \textbf{Final project} & \textbf{40\%} \\
  \textbf{Class participation} & \textbf{20\%} \\
  \bottomrule
\end{tabular}

\subsection{Assignments --- 40\% of final grade}

Four graded submissions, weighted equally. Assignments are submitted via Canvas.

\textbf{Late work receives a zero.} Assignment answers are worked through in class the following
afternoon, roughly fifteen hours after the deadline, which makes it impossible to accept a
submission after the fact --- by then the answers are public. There is no partial-credit sliding scale and no grace period. If something
comes up that will stop you meeting a deadline, email Joe as soon as you know, rather than after
the deadline has gone by.

\begin{tabular}{@{}llll@{}}
  \toprule
  \textbf{Assignment} & \textbf{Topic} & \textbf{Assigned} & \textbf{Due} \\
  \midrule
  Homework 1 & Logic & Sep 3 & Wed Sep 16, 11:59pm \\
  Homework 2 & Semantic networks \& search & Sep 24 & Wed Oct 7, 11:59pm \\
  Homework 3 & Expert system in CLIPS & Oct 15 & Wed Nov 4, 11:59pm \\
  Homework 4 & Bayesian networks & Oct 22 & Wed Nov 18, 11:59pm \\
  \bottomrule
\end{tabular}

\textbf{Everything is due at 11:59pm on a Wednesday}, the night before the class that reviews
it. The project proposal follows the same rule.

\subsection{Final project --- 40\% of final grade}

Either an AI-based research project (for doctoral students) or a knowledge-based system
implementation (for master's, professional, and other students), addressing a biomedical problem
of your choice. Full details below, and discussed in class. The 40\% breaks down as:

\begin{tabular}{@{}lrl@{}}
  \toprule
  \textbf{Component} & \textbf{Weight} & \textbf{Due} \\
  \midrule
  Proposal & 5\% & Wed Oct 21, 11:59pm \\
  Implementation and code & 10\% & Dec 11, 11:59pm \\
  Written report & 17\% & Dec 11, 11:59pm \\
  Recorded presentation & 8\% & Dec 11, 11:59pm \\
  \bottomrule
\end{tabular}

\subsection{Class participation --- 20\% of final grade}

\begin{tabular}{@{}lr@{}}
  \toprule
  \textbf{Component} & \textbf{Weight} \\
  \midrule
  Journal club presentation & 10\% \\
  Journal club discussion & 4\% \\
  In-class exercises and preparation & 6\% \\
  \bottomrule
\end{tabular}

The journal club presentation is worth as much as one homework assignment, which is roughly what
it costs to prepare. Discussion credit is for the other ten weeks --- everyone reads the paper,
not only the pair presenting it. In-class exercise credit is for engagement and completion, not
for correct answers.

Attendance is expected rather than credited. There are no points for showing up, so it is not
listed above. Each unexcused absence deducts 5\% from your final number grade for the course,
separately from the categories above.

There is no form to fill in and no list of approved reasons. Email Joe as soon as something
arises --- illness, a conflict, a family situation, anything --- and it will be treated as
excused. What turns an absence into an unexcused one is not the reason behind it but the silence:
absences become unexcused when you do not get in touch. The same goes for questions about
anything else in this syllabus. Ask early.

% -------------------------------------------------------- journal club ----
\section{Journal club}

Each week from Week 3 onward, we dedicate roughly 25 minutes to discussing one journal paper on
that week's topic. \textbf{This is half of what the course ran in previous years} --- one paper
per week instead of two --- so that we can spend the reclaimed time writing code.

A \textbf{pair} of students is assigned to each paper. One presents background and methods, the
other results and critique, and both lead the discussion afterward. With eleven papers and
roughly twenty students, everyone presents exactly once.

Paper presentations should include an overview of the paper (background, methods, results, and
the authors' conclusions), followed by a set of discussion questions posed to the rest of the
class. Presenters then lead the class in discussion and critique.

All students are expected to have read the paper and to participate actively in the discussion.
All papers are available in PDF format through the course website. The subject matter of each
paper is briefly discussed in class two weeks before the presentation date, to allow students
time to prepare.

\subsection{LLM-driven critique}

As part of each presentation, presenters must run the paper past an LLM chatbot of their choice,
ask it to critique the work, and report back on: what it flagged, which of those criticisms are
actually correct, and which are confidently wrong. The third question is the important one.

% -------------------------------------------------------- final project ----
\section{Final project}

The final project ties together multiple topics covered in this course, with a specific focus on
one or more sources or types of biomedical data. There are two types of final project, selected
based on your degree program or registration status.

\subsection{Research project (doctoral students; others may opt in)}

Students use the concepts taught over the semester to complete an original research project
addressing an area of interest or need in medicine. Examples include designing and evaluating a
new clinical decision support algorithm, developing an ontology or knowledge base to support a
specific domain along with use cases, or applying Bayesian networks to accomplish a novel task on
a biomedical dataset. All projects should be applied to either a real or synthetic dataset, and
should be cleared with Joe before work begins.

Note that \textbf{deductive inference must be used} in this project. A project that exclusively
or primarily focuses on machine learning is not acceptable.

\subsection{Implementation project (master's and professional students)}

Students propose and implement their own software tool to address a target biomedical task ---
for example, an expert system that makes deductions toward a clinical task, such as a decision
support system, an alert or reminder system, or one that proposes a diagnosis.

Students conduct a literature search targeting a specific biomedical task or problem, identify
relevant data, databases, experts, and resources for developing the tool, and complete an initial
implementation. The tool does not have to be production-ready, but should demonstrate the
intended functionality. This tool should be implemented in CLIPS (e.g., via clipspy).

\subsection{Proposal}

A one-page proposal is due \textbf{Wednesday, October 21, 2026, 11:59pm}, and is worth 5\% of the final
grade. It states the biomedical problem, where deductive inference enters, the data or resources
you plan to use, and which of the two project types you are doing. This is how your topic gets
cleared before you start building. It falls the week after CLIPS is covered in class, so
implementation students propose a tool having already worked with it.

\subsection{Final report}

Both project types involve a report with the same sections. The report should be an 8--10 page
description of the project, including:

\begin{itemize}
  \item \textbf{Background} --- synthesize relevant literature and the need for AI-based
        approaches to the problem. Include descriptions of any existing or previous deductive AI
        solutions in the area.
  \item \textbf{Materials and Methods} --- data, databases, and other resources used by and for
        the creation of the system.
  \item \textbf{Design} --- the high-level design of the proposed tool and how it integrates
        multiple topics covered in this course.
  \item \textbf{Demonstration} --- application of the tool to your task of interest and the
        results. If the tool is not completely implemented, or does not function as intended,
        discuss the current stage of implementation and your troubleshooting attempts.
  \item \textbf{Discussion and Conclusions} --- challenges you overcame, advantages and
        disadvantages of the tool, and how it could be improved or applied more broadly.
  \item \textbf{References} --- relevant citations. No specific format is required, but
        consistency and completeness both matter. Zotero or similar citation management software
        is recommended.
\end{itemize}

The report must be typeset in \LaTeX{} using the
\href{https://www.overleaf.com/latex/templates/aaai-press-latex-template/jymjdgdpdmxp}{AAAI
template}.

\textbf{Code submission.} Everyone additionally submits the code accompanying their project,
either (preferred) in a GitHub repository linked from the text of the paper, or as a documented,
well-structured directory alongside the report, compressed into a ZIP file.

\subsection{Oral presentation}

A 10-minute presentation structured around the format of the written report. Slides will be
helpful. The presentation is recorded as a video and uploaded to the Box folder linked from the
course website.

Your slides should include a brief introduction to the biomedical problem, how you approached
solving it using a deductive reasoning approach, and images (or a live demonstration) showing
your program's input and output. Research projects should also briefly describe the evaluation
approach and key takeaways.

Please stick to the 10-minute limit. Since you can edit and re-record, there is no reason to run
over, and points may be deducted for videos that run long.

\textbf{Both the report and the presentation video are due December 11, 2026.}

% ------------------------------------------------------ course materials ---
\section{Course materials}

No textbooks are strictly necessary, but the following are recommended:

\begin{itemize}
  \item Cawsey A. \emph{The Essence of Artificial Intelligence.} Pearson, 2nd ed., 2010.
  \item Russell S and Norvig P. \emph{Artificial Intelligence: A Modern Approach.} Pearson,
        4th ed., 2020.
\end{itemize}

Both may be found affordably online.

% ------------------------------------------------------ academic honesty ---
\section{Academic honesty}

All work submitted for credit is expected to be your own work. In the preparation of all papers
and other written work, you should always take great care to distinguish your own ideas and
knowledge from information derived from other sources. The term ``sources'' includes not only
published primary and secondary material, but also information and opinions gained directly from
other people. The responsibility for learning the proper forms of citation lies with you. You
must acknowledge any collaboration and its extent in all submitted work. You are expected to
follow Penn's standards of academic integrity as found in the
\href{https://catalog.upenn.edu/pennbook/code-of-academic-integrity}{Code of Academic Integrity}.

\subsection{Violations}

All violations of academic honesty and integrity are taken very seriously and will result in
disciplinary action. Depending on the seriousness of the violation, action will be taken at the
discretion of the course director, and may include reduction of grades, additional assignments, a
failing grade for the course, and/or reporting the incident to the University for further
evaluation.

% ------------------------------------------------------------ AI policy ----
\section{Statement on the use of artificial intelligence in assignments}

Although this is a course on AI, all work completed as part of the assignments, journal club, and
final project must be completed on your own volition, without the assistance of generative AI
tools (or similar), unless otherwise instructed by the course director. Any and all violations of
this policy will be treated similarly to other forms of plagiarism.

\textbf{Two exceptions:}

\begin{enumerate}
  \item The \textbf{LLM-driven critique} portion of each journal club, in which you are expected
        to use the LLM chatbot of your choice to assist in critiquing the article.
  \item The \textbf{in-class coding exercises}, which earn credit for engagement rather than for
        correct answers. Use whatever tools help you learn there --- though you will get more out
        of them by typing the code yourself.
\end{enumerate}

The homework assignments and the final project --- proposal, code, report, and presentation ---
are covered by the policy without exception.

Please direct any and all questions to Joe --- if you have any doubts, it is better to ask in
advance.

% -------------------------------------------------------- accommodations ---
\section{Students with disabilities}

The University of Pennsylvania provides reasonable accommodations to students with disabilities
who have self-identified and been approved by the office of Student Disabilities Services (SDS).
Please make an appointment to meet with me as soon as possible in order to discuss your
accommodations and your needs. If you have not yet contacted SDS and would like to request
accommodations or have questions, you can make an appointment by calling (215) 573-9235. The SDS
office is located in the Weingarten Learning Resources Center at Stouffer Commons, 3702 Spruce
Street, Suite 300. All SDS services are free and confidential. Please visit the
\href{https://www.vpul.upenn.edu/lrc/sds/}{SDS website} for more information.

% ------------------------------------------------------------- licensing ---
\section{Course materials: reuse and redistribution}

Course materials fall into two groups, and they are not treated the same way.

\textbf{Public materials --- reusable.} Everything in the public course repository on GitHub ---
the in-class notebooks, the course website, and the structure and organization of the course
itself --- is released under a
\href{https://creativecommons.org/licenses/by-nc-sa/4.0/}{Creative Commons
Attribution-NonCommercial-ShareAlike 4.0} licence. You may share and adapt this material,
including after the course ends, provided you give attribution, do not use it commercially, and
license anything you build from it on the same terms. Instructors elsewhere are welcome to it.

\textbf{Anything behind a Penn login --- not redistributable.} Material hosted on Canvas or Penn
Box, including the lecture slides, the journal club papers, the assignment files, and the class
recordings, is available to enrolled students only and must not be redistributed. This is not a
formality: much of that material is copyrighted by publishers and others, and the course director
holds neither the copyright nor a licence to pass it on. Posting a journal club PDF to a public
repository or a file-sharing site is a copyright violation regardless of the course's own
licensing.

If you are unsure which group something belongs to, the rule of thumb is where you got it: if
PennKey was required, treat it as not redistributable.

% --------------------------------------------------- program policies ------
\section{University and program policies}

This syllabus covers the policies specific to BMIN 5200. Policies that are not addressed here are
governed by the \href{\handbookurl}{IBI Master's and Certificate Programs Student and Faculty
Handbook} (2026--2027 edition), which is the reference for anything this document does not say.
Among other things, it covers the Code of Academic Integrity and Code of Student Conduct, the
University's Equal Opportunity and Nondiscrimination Statement, the Sexual Misconduct Policy and
associated resource offices and complaint procedures, student and academic grievance procedures,
registration, leave of absence, audit, academic standing and probation, the program's
standardized grading scale, course evaluations, the policy on secular and religious holidays,
suspension of normal operations and University closures, guidance on academic issues for students
with disabilities, and University resources including Counseling and Psychological Services
(CAPS) and the Weingarten Learning Resources Center.

The handbook is available at:\\
\href{\handbookurl}{\texttt{\handbookurl}}

Note that the handbook governs the Biomedical Informatics master's and certificate programs (MCI,
MSBMI, and the BMI Certificate). Students enrolled through a doctoral program or another school
should follow their own graduate group's policies where the two differ; the course-specific
policies in this syllabus apply to everyone.

% ------------------------------------------------------------- schedule ----
\section{Schedule at a glance}

Journal club papers are deliberately one week behind the lecture they cover: the paper discussed
in a given week is on the previous week's topic, so that presenters have already been taught the
material they are presenting on. Weeks 1 and 2 have no paper.

The Due column marks the meeting an item is due \emph{for}. The deadline itself is 11:59pm the
night before --- the Wednesday.

\begin{longtable}{@{}cl>{\raggedright\arraybackslash}p{5.1cm}>{\raggedright\arraybackslash}p{3.3cm}l@{}}
  \toprule
  \textbf{\#} & \textbf{Date} & \textbf{Topic} & \textbf{Journal club} & \textbf{Due} \\
  \midrule
  \endfirsthead
  \toprule
  \textbf{\#} & \textbf{Date} & \textbf{Topic} & \textbf{Journal club} & \textbf{Due} \\
  \midrule
  \endhead
  1 & Aug 27 & Course intro; history of AI & --- & \\
  2 & Sep 3 & Knowledge representation \& logic & --- & \\
  3 & Sep 10 & Semantic networks, frames, ontologies & Babalou et al. & \\
  4 & Sep 17 & Heuristic, local \& population-based search & Wolpert \& Macready & \textbf{HW 1} \\
  5 & Sep 24 & Biologically-inspired search & Nagarajan \& Babu & \\
  --- & Oct 1 & \emph{No class --- Fall Break} & & \\
  6 & Oct 8 & Rules \& knowledge-based systems & Shortliffe et al. & \textbf{HW 2} \\
  7 & Oct 15 & Building an expert system: CLIPS / clipspy & Michalowski et al. & \\
  8 & Oct 22 & Bayesian networks; state machines & Leclerc et al. & \textbf{Project proposal} \\
  9 & Oct 29 & Information theory \& machine learning & Shen et al. & \\
  10 & Nov 5 & Deep learning \& large language models & Vaswani et al. & \textbf{HW 3} \\
  11 & Nov 12 & Explainable AI & Behrad et al. & \\
  12 & Nov 19 & Bias \& fairness in AI & Pfohl et al. & \textbf{HW 4} \\
  --- & Nov 26 & \emph{No class --- Thanksgiving} & & \\
  13 & Dec 3 & Agentic AI & Thirunavukarasu et al. & \\
  --- & Dec 11 & & & \textbf{Final report + video} \\
  \bottomrule
\end{longtable}

\end{document}
